\section{Measures and Scoring}
\label{sec:scoring}
\subsection{Task Performance}
\paragraph{Active completion time.} The primary outcome is elapsed active-search time in seconds from object reveal to the correct right-trigger selection, excluding the union of recorded pause intervals overlapping that period. Section~\ref{ch:measurement} specifies the clock equation and handling of overlapping pauses. Lower values indicate faster completion. Readiness, opening guidance, post-selection feedback, questionnaires and between-trial breaks do not contribute to this outcome. An unfinished trial retains its observed exposure and outcome, but does not receive a successful completion time. The measurement definition also includes the full interval from opening-cue onset to selection to describe the complete instructed task duration.

\paragraph{Target contact and selection errors.} Time to first target gaze is the active-clock time of the first valid, recorded gaze-ray intersection with the target, and time to first controller contact uses the first tracked right-controller ray--target intersection. Last target-gaze contact time denotes the last observed valid target intersection during the trial, not a fixation duration or a spoken confirmation. These event times remain distinct from trigger selection and from the larger gaze-vicinity region used by the guidance system. Incorrect selections are counted as trigger presses that select non-target objects. If no valid target contact is observed, its latency remains missing rather than being set to zero, and all contact times retain their source and sampling resolution.

\subsection{Starting Geometry and Search Behavior}
\paragraph{Starting angle, height and distance.} The pre-cue measurement contract requires a saved reference origin, readiness-cross position and target position. Starting theta ($\theta_0$) is the signed angle about world up between the horizontal origin-to-cross and origin-to-target vectors; its magnitude $|\theta_0|$ quantifies horizontal turn demand, and its sign preserves direction. Degenerate horizontal vectors are marked unavailable. The primary analysis adjusts for $|\theta_0|$ rather than dividing completion time by degrees. Nominal wall-angle assignment, exact pre-cue theta and post-reveal head-to-target error remain separate variables. The exact pre-cue snapshot is not yet implemented in the aligned study build, and existing nominal-theta fields are not substitutes for it; the snapshot must be implemented and verified before the study release is frozen.

Signed target-height difference is $\Delta h=y_{\mathrm{target}}-y_{\mathrm{cross}}$, in metres; positive values indicate above-cross targets. Euclidean target distance is the three-dimensional distance between target center and the first valid tracked head position after reveal, with that observation's timestamp retained. All positions used in a calculation share a recorded coordinate frame.

\paragraph{Head movement and zone coverage.} Gross head translation sums distances between adjacent valid head positions, in metres, and gross head rotation sums shortest quaternion angular displacements between adjacent valid orientations, in degrees. Fixed-axis component totals are retained separately and are not mislabeled as Euler yaw, pitch and roll. Invalid intervals, pauses and excessive sample gaps break accumulation; the analysis does not interpolate across them. Covered time and gap counts accompany totals. These measures describe head movement, not independently tracked chair rotation.

Distinct walls visited, wall transitions and revisits summarize search coverage using stable valid gaze observations and a versioned zone-assignment rule. Transitions across tracking gaps are not inferred.

\paragraph{Instructions and guidance timing.} The delivered active-search assistance total counts playback starts for directional guidance, gaze/pointing acknowledgments, and selection reminders, including delivered repeats; these categories are also reported separately, as specified in Section~\ref{ch:measurement}. Requested repeats, delivered repeats and interrupted utterances are retained separately; model requests, prefetched audio and success speech are not counted as delivered search instructions. Opening cues and success feedback retain separate event labels. Timing analysis distinguishes evidence qualification, decision-request latency, model deferral, scheduling delay and audio onset. Gaze confirmations and controller confirmations are evaluated separately against their corresponding gates; the controller gate includes both its hold and the required nearby valid gaze.

\subsection{Subjective Measures}
\paragraph{Workload.} Six NASA-TLX dimensions---mental, physical and temporal demand, performance, effort and frustration---are rated separately for each voice block during the final survey, after both blocks. The study uses a documented 0--10 adaptation without pairwise weighting. Demand, effort and frustration run from 0 = Very low to 10 = Very high. The Performance prompt asks how successful the participant was and runs from 0 = Failure to 10 = Perfect. This favorable-direction response is a deliberate local adaptation of the opposite orientation in the \href{https://human-factors.arc.nasa.gov/groups/TLX/downloads/TLXScale.pdf}{NASA rating form}. All six raw responses are retained; higher raw Performance values indicate better perceived performance. Before calculating the workload mean, only Performance is reverse-scored once as $10-r_{\mathrm{Performance}}$. The unweighted mean uses the other five raw ratings and this reversed Performance value, and is calculated only when all six responses are present. Higher mean scores therefore indicate greater workload or poorer perceived performance. The source instrument is \href{https://doi.org/10.1016/S0166-4115(08)62386-9}{Hart and Staveland's NASA-TLX}.

\paragraph{Motivation and activity experience.} The \href{https://selfdeterminationtheory.org/intrinsic-motivation-inventory/}{Intrinsic Motivation Inventory} contributes interest/enjoyment (7 items), effort/importance (5), pressure/tension (5), and value/usefulness (7). Responses range from 1 = Not at all true to 7 = Very true. Published reverse keys are applied before averaging each complete subscale. Interest/enjoyment represents intrinsic motivation; the four subscales are not combined into an overall IMI score. Activity wording refers to the completed search block, with task-specific adaptations documented.

\paragraph{Agent perception.} Twenty-one items from \href{https://doi.org/10.1145/3651288}{Guo et al. (2024), Appendix A, Table A1} assess co-presence (Q1--6), perceived intelligence (Q13--18), appealing/eerie/familiar impressions (Q20--22), and perceived anthropomorphism (Q37--42). ``Virtual character'' is adapted to ``voice agent.'' Responses range from 1 = Very strongly disagree to 7 = Very strongly agree. The three six-item constructs have separate complete-scale means. Q20--22 are analyzed individually, without a combined eeriness score.

\paragraph{Voice experience and reliance.} Twelve custom 1--7 agreement items assess voice similarity, responsiveness, helpfulness, distraction, familiarity, trust, comfort, uncanniness, reliance, naturalness, willingness to follow and motivation to act without guidance. Voice similarity serves as a manipulation check. Voice liking and retrospective percentage of guidance use each use 0--100 responses with separate unable-to-answer options. These 14 numeric items are analyzed separately; they are not a validated composite scale. Liking during voice audition remains separate from condition-specific liking collected at the end.

\paragraph{Debrief.} Open debrief questions explore perceived performance, guidance experience, own-voice comfort, possible use cases and voice preference; broad experience questions precede specific eeriness or inner-dialogue probes.

Each condition-specific set in the final survey contains 65 numeric responses (130 across both sets): six workload ratings, 24 IMI items, 21 Guo items and 14 custom ratings. Technical checks and optional text are additional. Questionnaire observations have participant-by-block scope even when copied into per-trial summary files; they are not treated as 40 independent responses. Missing answers are not replaced with zero.


Table~\ref{tab:battery} summarizes the condition-specific battery. Counts refer to numeric items in one block-referenced set, not independent scales. Exact wording and reverse-key positions are preserved in Appendices~\ref{app:instruments} and~\ref{app:data}.
\begin{table}[htbp]
\centering
\caption[Condition-specific questionnaire battery]{Proposed final questionnaire battery for each referenced voice block. These ratings are administered after both blocks; setup voice ratings are separate.}
\label{tab:battery}
\small
\tagpdfsetup{table/header-rows={1}}
\begin{tabularx}{\linewidth}{@{}p{0.25\linewidth}p{0.09\linewidth}X@{}}
\toprule
Measure & Items & Scoring and interpretation \\\midrule
NASA-TLX adaptation & 6 & Six 0--10 dimensions; complete unweighted mean after reversing only raw Performance as $10-r$. Preserve raw responses and explicit bounds. \\
IMI & 24 & Separate interest/enjoyment (7), effort/importance (5), pressure/tension (5), and value/usefulness (7) means. No overall total. \\
Agent perception & 21 & Three six-item means: co-presence, intelligence, anthropomorphism. Appealing, eerie, and familiar remain three single items. \\
Custom voice and assistance & 14 & Individual ratings of voice and guidance experience. No unvalidated combined score. \\\midrule
Total per block reference & 65 & 130 numeric ratings across two references, before technical checks and optional text. \\\bottomrule
\end{tabularx}
\end{table}

\subsection{Measurement Quality}
Tracking validity, observation age, data coverage, gaps and spatial-validation error accompany gaze-derived measures. Ray intersections and vicinity hits are not labeled fixations, and render-frame observations do not establish native eye-tracker sampling frequency. Pupil-based workload inference and fixation/saccade outcomes require independent validation before inclusion. Trial performance and subjective summaries remain linked to separate raw gaze, motion and event records through participant, session, block and trial identifiers. Scale metadata distinguishes the present 0--10 workload form from any earlier exports using another range.

